%% ============================================================
%% AutoEval-Modeling · 自动生成论文主模板
%% 模板引擎: Jinja2 (block/variable 标记: <%  %> / <<  >>)
%% ============================================================
\documentclass[12pt, a4paper]{article}

%% -------------------- 宏包 --------------------
\usepackage[UTF8]{ctex}                 % 中文支持
\usepackage{amsmath, amssymb, amsthm}   % 数学
\usepackage{booktabs}                   % 三线表
\usepackage{graphicx}                   % 图片
\usepackage{float}                      % 浮动控制
\usepackage{geometry}                   % 页面设置
\usepackage{hyperref}                   % 超链接
\usepackage{caption}                    % 图表标题
\usepackage{subcaption}                 % 子图
\usepackage{longtable}                  % 长表格
\usepackage{array}                      % 表格列格式
\usepackage{multirow}                   % 合并行
\usepackage{xcolor}                     % 颜色
\usepackage{enumitem}                   % 列表样式
\usepackage{fancyhdr}                   % 页眉页脚
\usepackage{appendix}                   % 附录
\usepackage{listings}                   % 代码
\usepackage{xcolor}

\geometry{left=2.5cm, right=2.5cm, top=2.5cm, bottom=2.5cm}

%% -------------------- 代码样式 --------------------
\lstset{
    language=Python,
    basicstyle=\small\ttfamily,
    keywordstyle=\color{blue}\bfseries,
    commentstyle=\color{gray},
    stringstyle=\color{red},
    numbers=left,
    numberstyle=\tiny\color{gray},
    frame=single,
    breaklines=true,
    captionpos=b
}

%% -------------------- 页眉页脚 --------------------
\pagestyle{fancy}
\fancyhf{}
\fancyhead[C]{<< title >>}
\fancyfoot[C]{\thepage}
\renewcommand{\headrulewidth}{0.4pt}

%% -------------------- 标题 --------------------
\title{\textbf{<< title >>}}
\author{<< author | default("AutoEval-Modeling 自动生成") >>}
\date{<< date | default("\\today") >>}

\begin{document}

\maketitle

%% ==================== 摘要 ====================
\begin{abstract}
<< abstract_content >>

\textbf{关键词：}<< keywords | join("；") >>
\end{abstract}

\tableofcontents
\newpage

%% ==================== 一、问题重述 ====================
\section{问题重述}

\subsection{问题背景}
<< problem_background >>

\subsection{问题要求}
<< problem_requirements >>

%% ==================== 二、问题分析 ====================
\section{问题分析}
<< problem_analysis >>

%% ==================== 三、模型假设 ====================
\section{模型假设}
\begin{enumerate}[label=(\arabic*)]
<% for assumption in assumptions %>
    \item << assumption >>
<% endfor %>
\end{enumerate}

%% ==================== 四、符号说明 ====================
\section{符号说明}
\begin{table}[H]
    \centering
    \caption{主要符号说明}
    \label{tab:symbols}
    \begin{tabular}{cl}
        \toprule
        \textbf{符号} & \textbf{含义} \\
        \midrule
<% for sym in symbols %>
        $<< sym.notation >>$ & << sym.description >> \\
<% endfor %>
        \bottomrule
    \end{tabular}
\end{table}

%% ==================== 五、数据预处理 ====================
\section{数据预处理}
<< preprocessing_content >>

%% ==================== 六、模型建立与求解 ====================
\section{模型建立与求解}

<% for section in model_sections %>
<< section >>
<% endfor %>

%% ==================== 七、灵敏度分析 ====================
<% if sensitivity_content %>
\section{灵敏度分析}
<< sensitivity_content >>
<% endif %>

%% ==================== 八、模型评价与改进 ====================
\section{模型评价与改进}

\subsection{模型优点}
\begin{enumerate}[label=(\arabic*)]
<% for pro in model_pros %>
    \item << pro >>
<% endfor %>
\end{enumerate}

\subsection{模型缺点与改进方向}
\begin{enumerate}[label=(\arabic*)]
<% for con in model_cons %>
    \item << con >>
<% endfor %>
\end{enumerate}

%% ==================== 参考文献 ====================
\newpage
\begin{thebibliography}{99}
<% for ref in references %>
    \bibitem{<< ref.key >>} << ref.text >>
<% endfor %>
\end{thebibliography}

%% ==================== 附录 ====================
<% if appendix_content %>
\newpage
\begin{appendices}
\section{核心代码}
<< appendix_content >>
\end{appendices}
<% endif %>

\end{document}